\uuid{k5Sm}
\exo7id{159}
\titre{exo7 159}
\auteur{cousquer}
\organisation{exo7}
\datecreate{2003-10-01}
\isIndication{false}
\isCorrection{false}
\chapitre{Logique, ensemble, raisonnement}
\sousChapitre{Récurrence}
\cpge{mpsi-mp2i.raisonnement-vocabulaire-ensembliste.a, pcsi.raisonnement-vocabulaire-ensembliste.a, ptsi.raisonnement-vocabulaire-ensembliste.a, mpsi-mp2i.calcul-algebrique-trigonometrie.a, pcsi.calcul-algebrique-trigonometrie.a, ptsi.calcul-algebrique-trigonometrie.a, bcpst1.logique-ensembles-raisonnement.c}
\module{Algèbre}
\niveau{L1}
\difficulte{}

\contenu{
\question{
Pour tout entier naturel $n$, on pose 
$$ S_{n} = 1\cdot 2 + 2 \cdot 3+\cdots + (n-1)\cdot n$$
Démontrer que l'on a 
$$S_{n} = \frac{1}{3} n(n-1)(n+1)$$
}
}
